% Template for post or presentation abstracts submitted to ILAS 2022
% 1. Fill in the presenter's information (name, affiliation, email
% address) and talk information (title of presentation, abstract).
% 2. Compile to make sure there are no errors.
% 3. Save the .tex file for your abstract with a distinctive name,
% ideally "FamilynameGivenname.tex".
% 4. Upload your .tex file to https://clr.ie/129557
%       
% * Please keep your LaTeX commands simple, and avoid the use of
%    user-defined macros.
% * Include details of co-authors and funding acknowledgements after the abstract.
% * Upload only the LaTeX file (with .tex extension).
% * Questions: Contact Niall Madden (Niall.Madden@NUIGalway.ie)

\documentclass[a4paper]{amsart} 
\usepackage{amssymb} 
\usepackage{hyperref}
\pagestyle{empty}
\date{} 

%% IGNORE THE NEXT 4 LINES
\makeatletter % 
\def\labelenumi{\theenumi.}
\def\theenumi{\@arabic\c@enumi}
\makeatother
%% STOP IGNORING NOW

\begin{document}

\title{Null-Space Projects for Intermediate Students: Tomography, Cryptography, and More}
\author{Tom Asaki} %Presenting author only
\address{Washington State University} % Affiliation of presenting author
\email{tasaki@wsu.edu} % Email address of presenting author only

\maketitle

For many students, a non-trivial nullspace is a ``necessary evil'' associated with a non-injective linear transformation, $T:\mathbb{R}^n\rightarrow\mathbb{R}^m$.  They understand that such transformations are not invertible, but strategies, such as least squares methods and pseudo-inversion via singular value decomposition, are designed to avoid nullspace contributions using approximate inverse transformations, $P:\mathbb{R}^m\rightarrow\mathbb{R}^n$.  And, theorems even demonstrate that these methods represent best strategies, in that, recovered domain-space vectors are optimal.  That is, if $T(x)=b$ and $\hat{x}=P(T(x))$, then $\|\hat{x}-x\|_2$ is minimized over all possible linear transformations $P$.  In this work we describe improved pseudo-inversion methods that directly incorporate nullspace vectors along with both actual and introduced prior knowledge.  We provide two examples suitable for curious students who wish more than nullspace avoidance.  These examples provide a springboard for additional project directions.

One project concept is improved tomographic reconstructions from radiographs.  The radiographic (approximate) linear transformation typically has a very large nullspace $n\gg m$.  Any pseudo-inverse transformation typically results in a reconstruction which is unacceptable without extensive post-processing.  We invite the student to use reasonable prior knowledge, such as accepting values from a given set, to find null vector contributions which enhance the result.  In many cases, dramatic improvement is realized, up to and including exact reconstructions.

A second project concept is the sending and deciphering of encrypted messages.  A message encrypted using a non-injective transformation is not simply recovered because of the loss of information.  However, with the introduction of an intertwined and simultaneously encrypted passphrase, the correct nullspace contribution can be recovered.  The transformation and encrypted message can be made public, while the private passphrase in known only to the sender and intended receiver.  


\end{document}


% Please leave the next bit alone  
%%% Local Variables: 
%%% mode: latex
%%% TeX-master: "../ILAS-2022"
%%% End:
